\section{Siły i momenty w przestrzeni wektorowej}

\subsubsection*{Iloczyn wektorowy}

\begin{equation}
    \Vec{c} = \Vec{a} \times \Vec{b} = \begin{vmatrix}
    \Vec{i} & \Vec{j} & \Vec{k} \\
    a_x & a_y & a_z \\
    b_x & b_y & b_z
    \end{vmatrix}  = \Vec{i}(a_y b_z - a_z b_y) - \Vec{j}(a_x b_z - a_z b_x) + \Vec{k}(a_x b_y - a_y b_x)
\end{equation}

\begin{equation}
    \Vec{c} = [a_y b_z - a_z b_y;\ a_z b_x - a_x b_z;\ a_x b_y - a_y b_x]
\end{equation}

\begin{equation}
    |\Vec{c}| = |\Vec{a}| |\Vec{b}| \sin\alpha
\end{equation}

\subsubsection*{Moment siły względem punktu}

Punkt, względem którego liczymy moment siły:

\begin{equation*}
    O(O_y;\ O_z;\ O_z)
\end{equation*}

Punkt zaczepienia siły:

\begin{equation*}
    A(A_y;\ A_z;\ A_z)
\end{equation*}

Ramię siły:

\begin{equation}
    \Vec{r} = \Vec{OA} = [A_x - O_x;\ A_y - O_y;\ A_z - O_z]
\end{equation}

Siła:

\begin{equation*}
    \Vec{F} = [F_x;\ F_y;\ F_z]
\end{equation*}

Moment siły względem punktu:

\begin{equation}
    \Vec{M}_{O,F} = \Vec{r} \times \Vec{F} = \begin{vmatrix}
    \Vec{i} & \Vec{j} & \Vec{k} \\
    r_x & r_y & r_z \\
    F_x & F_y & F_z
    \end{vmatrix}
\end{equation}

\subsubsection*{Redukcja układu do punktu}

Siła główna:

\begin{equation}
    \Vec{F}_0 = \sum \Vec{F}_i = \left[ \sum F_{i,x};\ \sum F_{i,y};\ \sum F_{i,z} \right]
\end{equation}

Moment główny:

\begin{equation}
    \Vec{M}_0 = \sum \Vec{r}_i \times \Vec{F}_i + \sum \Vec{M}_j
\end{equation}

gdzie $i$ - liczba sił skupionych w układzie, $j$ - liczba momentów skupionych w układzie.

\subsubsection*{Warunek równowagi statycznej}

\begin{equation}
    \begin{cases}
        \Vec{F}_0 = \Vec{0} \\
        \Vec{M}_0 = \Vec{0}
    \end{cases}
\end{equation}

\begin{equation}
    \begin{cases}
        F_{0,x} = 0 \\
        F_{0,y} = 0 \\
        F_{0,z} = 0 \\
        M_{0,x} = 0 \\
        M_{0,y} = 0 \\
        M_{0,z} = 0
    \end{cases}
\end{equation}